\roadmapSection{0}{COMMON FOUNDATIONS}{6--9 Weeks}

{\small\itshape\color{arligray} Do this first. Completely. Every advanced repair skill in this roadmap assumes fluency in everything below.}

% ══════════════════════════════════════════════════════════════════════════════
\roadmapSubSection{0.1}{Safety, ESD, and Workspace}{1 Week}

\roadmapHead{Essential Tools --- Buy Before Touching Any Board}
\begin{toolbadges}
\toolbadge{Multimeter UT61E+}
\toolbadge{Soldering Iron 936A+ / T12}
\toolbadge{Hot Air Station 858D / Quick 861DW}
\toolbadge{DC Power Supply RF4 RF-3005PRO}
\toolbadge{Anti-Static Mat (grounded)}
\toolbadge{Anti-Static Wrist Strap}
\toolbadge{ESD Tweezers SS-SA / Vetus}
\toolbadge{Desoldering Pump + Copper Wick}
\toolbadge{IPA 99\% Isopropyl Alcohol}
\toolbadge{No-Clean Flux Syringe (Amtech / NC-559)}
\toolbadge{Fume Extractor with HEPA filter}
\toolbadge{Fire-Safe Silicone Work Mat}
\end{toolbadges}

\roadmapHead{ESD --- Understand the Physics, Not Just the Rules}
\checkitem{Human body model (HBM): body capacitance ~150pF discharging through chip = instant gate oxide damage}{Chips fail silently under ESD --- intermittent faults weeks later, not immediate death. Hardest to diagnose.}
\checkitem{Machine model (MM) and charged device model (CDM): even grounded tools can discharge via CDM}{Any ungrounded metal touching a pin is a CDM risk event}
\checkitem{Anti-static mat + wrist strap: verify both connected to same ground point --- measure resistance mat-to-strap < 10M$\Omega$}{}
\checkitem{Ground the mat to mains earth via 1M$\Omega$ safety resistor --- limits fault current to safe level}{}
\checkitem{Never slide PCB across non-ESD surface --- friction tribocharges the board}{}
\checkitem{Store bare boards and ICs in anti-static bags or foam --- not bubble wrap, not paper}{}

\roadmapHead{Battery Safety --- Li-ion Failure Modes}
\checkitem{Thermal runaway: internal short $\to$ exothermic reaction $\to$ venting $\to$ fire. Irreversible once started}{Keep CO2 or dry sand nearby. Never use water on Li-ion fire.}
\checkitem{Discharge to below 20\% before any mechanical work on a phone or laptop}{}
\checkitem{Swollen battery = internal gas from electrolyte decomposition. Never compress, puncture, or charge}{Puncture in a sealed room = toxic gas (HF, CO). Work outdoors or with ventilation.}
\checkitem{Never short battery terminals --- even momentary: fuse wire melts, connector burns, PMIC can die}{}
\checkitem{Identify NTC thermistor on battery connector: 2-pin = bare cell (no protection), 3/4-pin = protected pack}{}
\checkitem{Dispose of swollen cells: submerge in salt water for 24h to discharge completely before disposal}{}

\roadmapHead{Workspace Setup}
\checkitem{Three zones: disassembly (mechanical) / soldering (heat) / testing (power on). Never mix zones.}{}
\checkitem{5000K LED lighting or natural light --- shadows hide lifted pads and hairline cracks}{}
\checkitem{Fume extractor: position nozzle 10--15cm from iron tip, pulling flux vapor away from face}{}
\checkitem{Label all screws and small parts with phone camera + sticky tape matrix before disassembly}{}

\newpage
% ══════════════════════════════════════════════════════════════════════════════
\roadmapSubSection{0.2}{Electronics Fundamentals --- Deep Level}{5--6 Weeks}

{\small\itshape\color{arligray} Every schematic, datasheet, and fault trace requires this. No shortcuts.}

\roadmapHead{Core Laws --- Apply, Not Just Memorize}
\checkitem{Ohm's Law: $V = IR$. Use it: if PP\_VBAT = 3.8V and charge IC sense resistor = 0.01$\Omega$, $I = V/R$}{}
\checkitem{KVL: sum of all voltages around any closed loop = 0. Use to find voltage drop across unknown element}{}
\checkitem{KCL: sum of currents into a node = sum out. Use to find current through unprobeable branch}{}
\checkitem{Power: $P = VI = I^2R = V^2/R$. A 0.01$\Omega$ sense resistor at 2A dissipates 0.04W --- safe. At 20A: 4W --- gets hot.}{}
\checkitem{Thevenin equivalent: any linear network = one voltage source + one series resistor. Use to model output impedance of power rails.}{}
\checkitem{Norton equivalent: same network = current source + parallel resistor. Use for current source analysis (charge IC output).}{}

\roadmapHead{Passive Components --- Test Each, Know Failure Modes}
\checkitem{Resistor: color code + DMM. Failure modes: open (most common after surge), value drift, intermittent (cold solder).}{In diode mode (board powered off): a shorted resistor reads 0V across it --- impossible unless bridged}
\checkitem{Capacitor: measure ESR with LCR meter (not just capacitance). High ESR = failing cap even if C value OK.}{Decoupling cap shorted to ground = rail pulled low. Most common cause of 0V on a power rail.}
\checkitem{Inductor: DC resistance in m$\Omega$ range. Open inductor = no output from buck converter. Shorted = overheating.}{}
\checkitem{Diode: forward voltage ~0.3V Schottky, ~0.6V silicon. Leaky diode = lower Vf than spec. Shorted = 0V both probes.}{}
\checkitem{Zener: reverse breakdown voltage = datasheet $V_Z$. Used for ESD clamping on USB, headphone, RF lines.}{}
\checkitem{Ferrite bead: acts as inductor at HF, resistor at RF. Shorted ferrite = fine at DC. Open ferrite = no power to downstream IC.}{}
\checkitem{NTC thermistor: $R$ drops with temperature. PMIC uses it to halt charging if battery temp exceeds threshold.}{}
\checkitem{Crystal oscillator: 32.768kHz RTC crystal, 19.2MHz / 38.4MHz system clock. Dead crystal = no boot, no time-keeping.}{}

\roadmapHead{Active Components --- Understand Operation Regions}
\checkitem{BJT (NPN/PNP): cutoff / active / saturation. In repair: used as switches (saturation). Base resistor sizing matters.}{}
\checkitem{MOSFET (N/P-channel): gate threshold $V_{GS(th)}$. Enhancement mode: off by default. Logic-level FETs: $V_{GS(th)}$ < 2V.}{Most common MOSFET fault: gate oxide punch-through = FET permanently on or off. Check $V_{GS}$ in circuit.}
\checkitem{Understand $R_{DS(on)}$: lower = less heat. Power MOSFETs in charger path: $R_{DS(on)}$ = 10--50m$\Omega$.}{}
\checkitem{Op-amp: understand virtual ground, inverting/non-inverting config. Used in charge IC current sense amplifiers.}{}

\roadmapHead{Power Architecture --- Shared Across Every Platform}
\checkitem{Buck (step-down) converter: $V_{out} = V_{in} \times D$ where $D$ = duty cycle. Switching freq = 500kHz--2MHz on phones.}{Inductor + capacitor form LC filter. Missing output = dead FET, missing gate drive, or open inductor.}
\checkitem{Boost (step-up) converter: $V_{out} = V_{in} / (1-D)$. Used for: display backlight (20--30V), USB output 5V from 3.7V battery.}{}
\checkitem{LDO (Low Dropout Regulator): linear, low noise. Dropout voltage = $V_{in} - V_{out(min)}$. Used for RF, audio, sensor rails.}{LDO gets hot if $V_{in} - V_{out}$ is large and current is high: $P = (V_{in}-V_{out}) \times I$}
\checkitem{Charge pump: capacitor-based voltage doubler/inverter. Used for: negative rail on audio codec, small bias voltages.}{}
\checkitem{PMIC (Power Management IC): integrates 5--15 buck/boost/LDO regulators + sequencer + fuel gauge + thermal protection}{On phone: one PMIC controls all major rails. Dead PMIC = completely dead device.}
\checkitem{Power sequencing: rails must power up in specific order (datasheet ``power-up sequence'' table). Wrong order = latch-up.}{Example iPhone: VBAT $\to$ PP1V8\_ALWAYS $\to$ PP3V0\_NAND $\to$ NAND responds $\to$ CPU fetches boot code}
\checkitem{Power rail naming convention: PP = power plane, voltage in name, function suffix. Example: PP1V8\_SDRAM, PP3V3\_USB}{}

\roadmapHead{Communication Protocols --- Read with Tools, Not Just Theory}
\checkitem{I2C: SDA + SCL, open-drain, pull-up to VDDIO. Probe with logic analyzer: verify ACK bit after each byte.}{Missing pull-up = SDA/SCL stuck low. Wrong pull-up value = signal too slow or too fast for IC.}
\checkitem{SPI: MOSI, MISO, SCK, CS (active low). Full duplex. Up to 100MHz. NAND flash, display ICs, PMIC config.}{No CS toggle = SPI device never selected. Wrong polarity (CPOL/CPHA) = scrambled data.}
\checkitem{UART: TX, RX, no clock. Match baud rate, 8N1 default. Boot logs always on UART TX at 115200.}{}
\checkitem{USB 2.0: D+/D- differential pair, 90$\Omega$ impedance. DP/DM short = USB dead. Open = enumeration fails.}{}
\checkitem{USB-C: CC1/CC2 lines negotiate role (host/device) and power delivery profile. Dead CC = no charging, no data.}{}
\checkitem{PCIe (NVMe): PERST\# must deassert, REFCLK must be present, PWRGT\# must enable. All three before link training.}{}
\checkitem{eMMC: CLK, CMD, DAT0--7. CLK stuck = storage dead. CMD high-Z = no command response. Use UFI Box to debug.}{}
\checkitem{UFS (phone/laptop): M-PHY differential pairs, MIPI UniPro framing. Oscilloscope needed for signal integrity.}{}

\roadmapHead{Schematic and Boardview Mastery}
\checkitem{Net names: identical labels on different pages = same copper. Never assume proximity = connection.}{}
\checkitem{Power flags: PWR\_FLAG in KiCad, circles in Altium. Means ``this net is a power plane''.}{}
\checkitem{Trace a charging path end-to-end: USB-C connector $\to$ CC controller $\to$ charge IC $\to$ battery connector $\to$ cell}{}
\checkitem{Identify all bypass caps on a rail: a shorted 100nF near the IC pulls the whole rail to ground}{}
\checkitem{Cross-reference schematic net to boardview component: click net in MaAnt, find physical location on board}{}
\checkitem{Identify test points (TP labels) and measure them first before probing IC pins directly}{}
\checkitem{Read IC application circuit in datasheet: compare to board schematic to spot missing or wrong components}{}

\newpage
% ══════════════════════════════════════════════════════════════════════════════
\roadmapSubSection{0.3}{Soldering and Micro-Soldering Mastery}{4--5 Weeks}

{\small\itshape\color{arligray} Daily practice on dead boards. 2 hours/day minimum. Cannot be read --- only built.}

\roadmapHead{Advanced Tools --- Buy at This Phase}
\begin{toolbadges}
\toolbadge{USB Microscope 45x / Yaxun AK12}
\toolbadge{Trinocular Stereo Microscope 7--45x (long-term)}
\toolbadge{DC Power Supply RF4 RF-3005PRO}
\toolbadge{Vacuum Separator Sunshine S-918K}
\toolbadge{Jumper Wire 0.02mm (Kaisi)}
\toolbadge{Flux Paste Amtech NC-559 / MF}
\toolbadge{Solder Paste Mechanic XG-Z40}
\toolbadge{BGA Reball Stencils (universal + platform-specific)}
\toolbadge{Low-Melt Solder (Chip Quik / Bismuth-Tin)}
\toolbadge{Preheat Platform (50--150°C bottom heat)}
\toolbadge{Precision Screwdrivers (Pentalobe / Torx / JIS)}
\toolbadge{10--15 Dead Practice Boards}
\end{toolbadges}

\roadmapHead{Iron Temperature and Tip Selection}
\checkitem{Lead-free (SAC305): 350--370°C. Leaded (Sn63Pb37): 300--320°C. Higher temp = faster oxidation of tip.}{}
\checkitem{Tip shapes: chisel for through-hole and drag soldering; conical for 0402 and fine pitch; knife for IC legs.}{}
\checkitem{Tin tip with fresh solder before and after every joint. Dull/oxidized tip = poor heat transfer = cold joints.}{}
\checkitem{Dwell time: 1--2 seconds per joint. Longer = pad delamination. Shorter = cold joint. Judge by solder flow.}{}
\checkitem{Flux before heat: always apply flux before touching iron to pad. Flux cleans oxide, improves wetting.}{}

\roadmapHead{SMD Component Work --- Build to These Standards}
\checkitem{0805: practice until removal and replacement is consistent with zero lifted pads --- 50+ repetitions}{}
\checkitem{0402: smallest common passive. Iron + hot air combo. Tweezers hold component while iron tacks one end.}{}
\checkitem{0201: phone-scale. Hot air only or iron with ultra-fine tip. Verify under 45x after every placement.}{}
\checkitem{SOT-23 (3-pin transistor / diode): hot air 350°C, 20--25L airflow. Rotate nozzle to avoid blowing off adjacents.}{}
\checkitem{QFN/DFN: exposed pad underneath. Hot air + flux paste + stencil. Center pad must solder for thermal + electrical.}{Verify center pad continuity with resistance measurement after rework}
\checkitem{TSSOP / SOP (IC legs): drag solder technique. Apply flux, load solder on tip, drag across all pins in one pass. Wick bridges.}{}

\roadmapHead{BGA --- Reball and Reflow Sequence}
\checkitem{Step 1: remove BGA using preheat platform (120°C) + hot air (380--400°C, 40--60L, circular motion)}{Never exceed 430°C --- PCB delaminates. Watch for slight movement before lifting.}
\checkitem{Step 2: clean pads with wick + flux. Inspect under microscope: all pads must be flat, no lifted pads.}{}
\checkitem{Step 3: apply flux to chip underside. Place stencil, apply solder paste, squeegee, remove stencil.}{}
\checkitem{Step 4: reflow paste in reflow oven or hot air. SAC305 profile: 150°C preheat 60s, 217°C liquidus, 250°C peak 10s.}{}
\checkitem{Step 5: clean flux residue with IPA + brush. Inspect all balls under microscope: uniform size, no bridges, no missing.}{}
\checkitem{Step 6: reinstall BGA. Align to pads using boardview as reference. Preheat + hot air reflow. No pressure during cooling.}{}
\checkitem{Step 7: verify electrically: power rail continuity, I2C/SPI bus ACK, or full boot depending on component.}{}

\roadmapHead{Micro-Soldering --- Jumper Wires and Pad Repair}
\checkitem{Lifted pad: scrape solder mask 3--5mm along the trace from the pad. Tin the exposed copper.}{}
\checkitem{0.02mm jumper wire: solder one end to tinned trace, route under adjacent components, solder other end to destination.}{}
\checkitem{Broken via: drill through from top with 0.3mm bit, insert wire, solder both ends.}{}
\checkitem{Rebuild a destroyed pad: clean area, apply UV solder mask, expose copper with scalpel, tin, attach component.}{}
\checkitem{Fine-pitch IC pin bridging: flux + solder wick to remove bridge. Never drag dry wick --- pre-tin it.}{}

\roadmapHead{Hot Air Station --- Precision Airflow Control}
\checkitem{Nozzle selection: 6mm round for small ICs; 10mm for medium; flat nozzle for connectors and flex cables.}{}
\checkitem{Prevent collateral damage: kapton tape + aluminum foil shields on all adjacent components within 15mm.}{}
\checkitem{Preheat from below: preheat platform at 100--130°C reduces thermal shock and time needed from hot air.}{}
\checkitem{Observe solder fillet on nearby passives as indicator: when they reflow, BGA solder is also liquid --- lift now.}{}
\checkitem{Cooling: allow board to cool naturally on heat-resistant mat. No forced air --- differential thermal contraction causes opens.}{}

\roadmapHead{PCB Layer and Trace Understanding}
\checkitem{Identify layer count from PCB edge cross-section or manufacturer spec. Phone PCBs: 8--12 layers typical.}{}
\checkitem{Blind via: connects top layer to inner layer but not through the board. Cannot be probed from bottom.}{}
\checkitem{Buried via: between two inner layers only. Invisible from outside. Identified from schematic only.}{}
\checkitem{Power plane (inner layer): large copper fill for VDD or GND. Low impedance path. Damaged by drill damage or delamination.}{}
\checkitem{Impedance-controlled traces: differential pairs (USB, PCIe, RF) have calculated width for 90/100$\Omega$. Do not add solder or flux.}{}
\checkitem{Identify trace resistance: $R = \rho L / A$. A 0.1mm wide, 1mm long copper trace has ~5m$\Omega$ resistance. Useful for fault current calc.}{}